% Part 2 of Day 1: Deep learning recap from a cost view.
%
% Topics of the syllabus: cost of each layer type; parameters, operations,
% and activation memory are three different budgets; efficient blocks:
% depthwise separable convolution and inverted residual. Exercise: calculate
% the parameters and the output size of three layers by hand.
%
% All sizes use 1 KB = 1024 bytes. The example network of this part is:
% input 96 x 96 x 3, three convolutions 3 x 3 with stride 2 and padding 1
% (16, 32, and 64 filters), global average pooling, one fully connected
% layer with 10 outputs.

\theorypart{Deep learning recap from a cost view}%
  {calculate the parameters and the output size of three layers}

% \cellgrid{x}{y}{cells}{cell size}{fill colour} draws a square grid with its
% lower left corner at (x,y). A frame cannot define a command with
% parameters, so the definition is here.
\providecommand{\cellgrid}[5]{%
  \fill[#5] (#1,#2) rectangle ++(#3*#4,#3*#4);
  \foreach \cellline in {0,...,#3} {
    \draw[black!70, line width=0.3pt] (#1+\cellline*#4,#2) -- ++(0,#3*#4);
    \draw[black!70, line width=0.3pt] (#1,#2+\cellline*#4) -- ++(#3*#4,0);
  }%
}

% -------------------------------------------------------------- three budgets
\begin{frame}{A model has three budgets}
  \begingroup\centering
  \begin{tikzpicture}[font=\small, >={Stealth[length=5pt]},
      m/.style={draw=cardinalred, line width=0.8pt, rounded corners=3pt,
                fill=boxgray, minimum width=3.0cm, minimum height=1.0cm,
                align=center},
      d/.style={draw=darkcyan, line width=0.8pt, rounded corners=3pt,
                fill=boxgray, minimum width=3.0cm, minimum height=1.0cm,
                align=center},
      q/.style={font=\footnotesize, text=coolgray, anchor=west, align=left}]
    \node[font=\footnotesize\bfseries, text=coolgray] at (0,1.0) {The model};
    \node[font=\footnotesize\bfseries, text=coolgray] at (4.2,1.0) {The device};
    \node[m] (p) at (0,0)    {Parameters};
    \node[m] (o) at (0,-1.4) {Operations};
    \node[m] (a) at (0,-2.8) {Activation\\memory};
    \node[d] (f) at (4.2,0)    {Flash};
    \node[d] (t) at (4.2,-1.4) {Time and\\energy};
    \node[d] (r) at (4.2,-2.8) {RAM};
    \draw[->] (p) -- (f);
    \draw[->] (o) -- (t);
    \draw[->] (a) -- (r);
    \node[q] at (5.9,0)    {Can the device store\\the model?};
    \node[q] at (5.9,-1.4) {Is the answer in time?\\Does the battery last?};
    \node[q] at (5.9,-2.8) {Do the tensors fit\\during inference?};
  \end{tikzpicture}\par\endgroup
  \vskip 0.6em
  \keypoint{The three budgets are different numbers. A model can meet two
    budgets and fail the third.}
\end{frame}

% ------------------------------------------------------------------- notation
\begin{frame}{Notation: tensors and the output size}
  \begin{columns}[c]
    \begin{column}{0.40\textwidth}
      \centering
      \begin{tikzpicture}[font=\footnotesize, line join=round]
        % a tensor as a box: front face, top face, side face
        \fill[boxgray] (0,0) rectangle (2.2,2.2);
        \draw[coolgray, line width=0.8pt] (0,0) rectangle (2.2,2.2);
        \draw[coolgray, line width=0.8pt, fill=boxgray]
          (0,2.2) -- (0.9,2.9) -- (3.1,2.9) -- (2.2,2.2) -- cycle;
        \draw[coolgray, line width=0.8pt, fill=boxgray]
          (2.2,0) -- (3.1,0.7) -- (3.1,2.9) -- (2.2,2.2) -- cycle;
        % kernel
        \draw[cardinalred, line width=1pt] (0.3,1.2) rectangle (1.0,1.9);
        \node[cardinalred] at (0.65,0.95) {$K \times K$};
        \node at (1.1,-0.3) {$W$};
        \node at (-0.3,1.1) {$H$};
        \node at (3.0,0.15) {$C$};
      \end{tikzpicture}
    \end{column}
    \begin{column}{0.58\textwidth}
      \begin{itemize}
        \item A feature map has $H \times W \times C$ values: height, width,
              channels.
        \item A kernel has the size $K \times K$.
        \item $S$ is the stride. $P$ is the padding.
      \end{itemize}
      \begin{block}{Output size of a convolution or a pooling layer}
        \[
          H_{\text{out}} = \left\lfloor \frac{H + 2P - K}{S} \right\rfloor + 1
        \]
        The same formula gives $W_{\text{out}}$.
      \end{block}
      Example: $H = 96$, $K = 3$, $S = 2$, $P = 1$.
      \[
        H_{\text{out}} = \lfloor 95 / 2 \rfloor + 1 = 48
      \]
    \end{column}
  \end{columns}
\end{frame}

% --------------------------------------------------------------- layer costs
\begin{frame}{Cost of a fully connected layer}
  \begin{columns}[T]
    \begin{column}{0.50\textwidth}
      A layer with $N_{\text{in}}$ inputs and $N_{\text{out}}$ outputs:
      \vskip 0.4em
      \renewcommand{\arraystretch}{1.3}
      \begin{tabular}{@{}ll@{}}
        \toprule
        Parameters  & $N_{\text{in}} \times N_{\text{out}} + N_{\text{out}}$ \\
        Operations  & $N_{\text{in}} \times N_{\text{out}}$ MACs \\
        Output size & $N_{\text{out}}$ values \\
        \bottomrule
      \end{tabular}
      \vskip 0.6em
      A MAC is one multiply-accumulate operation: $a \leftarrow a + w \times x$.
    \end{column}
    \begin{column}{0.46\textwidth}
      \begin{block}{Example: 784 inputs, 128 outputs}
        \begin{itemize}
          \item Weights: $784 \times 128 = 100\,352$
          \item Biases: 128
          \item Parameters: 100\,480
          \item Operations: 100\,352 MACs
          \item Output: 128 values
        \end{itemize}
      \end{block}
    \end{column}
  \end{columns}
  \vskip 0.6em
  \keypoint{Each weight is used one time. The number of operations is equal
    to the number of weights. The layer is expensive in flash.}
\end{frame}

\begin{frame}{Cost of a convolution layer}
  \begin{columns}[T]
    \begin{column}{0.54\textwidth}
      $C_{\text{out}}$ filters, each of size $K \times K \times
      C_{\text{in}}$:
      \vskip 0.4em
      \footnotesize
      \renewcommand{\arraystretch}{1.4}
      \begin{tabular}{@{}ll@{}}
        \toprule
        Parameters  & $K^2 \, C_{\text{in}} \, C_{\text{out}} + C_{\text{out}}$ \\
        Operations  & $K^2 \, C_{\text{in}} \, C_{\text{out}} \times
                      H_{\text{out}} \, W_{\text{out}}$ MACs \\
        Output size & $H_{\text{out}} \times W_{\text{out}} \times
                      C_{\text{out}}$ values \\
        \bottomrule
      \end{tabular}
    \end{column}
    \begin{column}{0.42\textwidth}
      \begin{block}{Example: $3 \times 3$, 32 to 64 channels, output
          $56 \times 56$}
        \begin{itemize}
          \item Weights: $9 \times 32 \times 64 = 18\,432$
          \item Parameters: 18\,496
          \item Operations: $18\,432 \times 3136 \approx 57.8$ million MACs
          \item Output: 200\,704 values
        \end{itemize}
      \end{block}
    \end{column}
  \end{columns}
  \vskip 0.6em
  \keypoint{The layer uses each weight again at each output position. It is
    small in flash, but expensive in operations and in RAM.}
\end{frame}

\begin{frame}{The two layer types have opposite costs}
  \renewcommand{\arraystretch}{1.3}
  \begin{tabular}{@{}L{0.30\textwidth}L{0.30\textwidth}L{0.30\textwidth}@{}}
    \toprule
    & \textbf{Fully connected} & \textbf{Convolution} \\
    \midrule
    The parameters grow with & the number of inputs & the kernel size and
      the channels \\
    A larger image gives     & more parameters & the same parameters, more
      operations \\
    Operations for each weight & 1 & $H_{\text{out}} \times W_{\text{out}}$ \\
    The main cost            & flash & operations and RAM \\
    \bottomrule
  \end{tabular}
  \vskip 0.8em
  \begin{itemize}
    \item Example: a fully connected layer from 1024 inputs to 1024 outputs
          has 1\,048\,576 weights.
    \item A convolution layer with a $3 \times 3$ kernel shares its weights
          across the image. Its parameters do not depend on the image size.
  \end{itemize}
\end{frame}

\begin{frame}{Layers with few parameters or no parameters}
  \small
  \renewcommand{\arraystretch}{1.35}
  \begin{tabular}{@{}L{0.22\textwidth}L{0.17\textwidth}L{0.52\textwidth}@{}}
    \toprule
    \textbf{Layer} & \textbf{Parameters} & \textbf{Effect on the cost} \\
    \midrule
    Activation (ReLU) & 0 & One comparison for each value. The size does not
      change. \\
    Pooling, stride 2 & 0 & The output has one quarter of the values. All
      later layers cost less. \\
    Global average pooling & 0 & $H \times W \times C$ values become $C$
      values. It replaces a large fully connected layer. \\
    Batch normalization & $2C$ & After training, the converter merges it
      into the convolution before it. Its cost at inference is then zero. \\
    \bottomrule
  \end{tabular}
  \vskip 0.8em
  \keypoint{A stride of 2 early in the network is the cheapest method to
    decrease the operations and the RAM of all later layers.}
\end{frame}

% ------------------------------------------------------------ example network
\begin{frame}{Example: a small network, layer by layer}
  Input $96 \times 96 \times 3$. Each convolution: $3 \times 3$, stride 2,
  padding 1.
  \vskip 0.5em
  \small
  \renewcommand{\arraystretch}{1.25}
  \begin{tabular}{@{}lrrrr@{}}
    \toprule
    \textbf{Layer} & \textbf{Output} & \textbf{Values} & \textbf{Parameters}
      & \textbf{MACs} \\
    \midrule
    Input                  & $96 \times 96 \times 3$  & 27\,648 & & \\
    Convolution 1, 16 filters & $48 \times 48 \times 16$ & 36\,864 & 448
      & 995\,328 \\
    Convolution 2, 32 filters & $24 \times 24 \times 32$ & 18\,432 & 4640
      & 2\,654\,208 \\
    Convolution 3, 64 filters & $12 \times 12 \times 64$ & 9216 & 18\,496
      & 2\,654\,208 \\
    Global average pooling & 64 & 64 & 0 & \\
    Fully connected        & 10 & 10 & 650 & 640 \\
    \midrule
    \textbf{Total}         &    &    & \textbf{24\,234}
      & \textbf{6\,304\,384} \\
    \bottomrule
  \end{tabular}
  \vskip 0.6em
  \normalsize
  Calculate one row with the formulas before you continue.
\end{frame}

\begin{frame}{Example: where the costs are}
  \begingroup\centering
  \begin{tikzpicture}[font=\footnotesize, x=0.92cm, y=0.027cm]
    % y unit: 1 = 1 percent. Three bars for each layer.
    \draw[coolgray] (0,0) -- (10.2,0);
    \foreach \v in {0,25,50,75,100} {
      \draw[coolgray!50] (0,\v) -- (10.2,\v);
      \node[left, coolgray] at (0,\v) {\v\,\%};
    }
    % convolution 1
    \fill[cardinalred!75] (0.4,0) rectangle (1.0,1.8);
    \fill[darkcyan!75]    (1.0,0) rectangle (1.6,15.8);
    \fill[treegreen!75]   (1.6,0) rectangle (2.2,100);
    \node[below] at (1.3,0) {Convolution 1};
    % convolution 2
    \fill[cardinalred!75] (2.9,0) rectangle (3.5,19.1);
    \fill[darkcyan!75]    (3.5,0) rectangle (4.1,42.1);
    \fill[treegreen!75]   (4.1,0) rectangle (4.7,85.7);
    \node[below] at (3.8,0) {Convolution 2};
    % convolution 3
    \fill[cardinalred!75] (5.4,0) rectangle (6.0,76.3);
    \fill[darkcyan!75]    (6.0,0) rectangle (6.6,42.1);
    \fill[treegreen!75]   (6.6,0) rectangle (7.2,42.9);
    \node[below] at (6.3,0) {Convolution 3};
    % pooling and fully connected
    \fill[cardinalred!75] (7.9,0) rectangle (8.5,2.7);
    \fill[darkcyan!75]    (8.5,0) rectangle (9.1,0.3);
    \fill[treegreen!75]   (9.1,0) rectangle (9.7,14.4);
    \node[below] at (8.8,0) {Pooling and FC};
    % legend
    \fill[cardinalred!75] (0.4,118) rectangle (0.7,126);
    \node[anchor=west] at (0.75,122) {parameters (share of the total)};
    \fill[darkcyan!75] (5.6,118) rectangle (5.9,126);
    \node[anchor=west] at (5.95,122) {MACs (share of the total)};
    \fill[treegreen!75] (0.4,106) rectangle (0.7,114);
    \node[anchor=west] at (0.75,110) {RAM: input and output of the layer
      (share of the maximum)};
  \end{tikzpicture}\par\endgroup
  \vskip 0.3em
  \begin{itemize}
    \item The first layer has 2\,\% of the parameters, but it needs the most
          RAM.
    \item The last convolution has 76\,\% of the parameters.
  \end{itemize}
  \keypoint{The early layers set the RAM budget. The late layers set the
    flash budget.}
\end{frame}

% -------------------------------------------------------------------- budget 1
\begin{frame}{Budget 1: parameters and flash}
  \begin{block}{Rule}
    Model size $\approx$ number of parameters $\times$ bytes for each
    parameter.
  \end{block}
  \begin{columns}[T]
    \begin{column}{0.40\textwidth}
      \begin{itemize}
        \item \code{float32}: 4 bytes
        \item \code{int8}: 1 byte (Day 4)
        \item The model file also holds the structure of the network. This
              part is small.
      \end{itemize}
    \end{column}
    \begin{column}{0.56\textwidth}
      \small
      \renewcommand{\arraystretch}{1.25}
      \setlength{\tabcolsep}{4pt}
      \begin{tabular}{@{}lrr@{}}
        \toprule
        \textbf{Model} & \textbf{Parameters} & \textbf{float32} \\
        \midrule
        Example network  & 24\,234     & 95 KB \\
        Keyword spotting & 26\,000     & $<$ 100 KB \\
        MobileNetV2      & 3.5 million & 13.4 MB \\
        ResNet-50        & 25.6 million & 98 MB \\
        \bottomrule
      \end{tabular}
    \end{column}
  \end{columns}
  \vskip 0.6em
  On a microcontroller, the flash holds the program and the model. The
  weights stay in the flash. The device does not copy them to the RAM.
  \source{Parameters of the three models: \emph{Machine Learning Systems},
    Volume I, chapters 2 and 6. Sizes with 1 MB = 1024 KB.}
\end{frame}

% -------------------------------------------------------------------- budget 2
\begin{frame}{Budget 2: operations, time, and energy}
  \begin{columns}[T]
    \begin{column}{0.50\textwidth}
      \begin{itemize}
        \item Count the MACs of each layer and add them.
        \item One MAC is two floating-point operations (FLOPs).
        \item Many papers write ``FLOPs'' for the number of MACs. Read the
              definition before you compare.
        \item More operations need more time and more energy on the same
              device.
      \end{itemize}
    \end{column}
    \begin{column}{0.46\textwidth}
      \begin{block}{The count is an estimate}
        \begin{itemize}
          \item The time also depends on memory access, on the kernel
                library, and on the data type.
          \item Example of the book: MobileNetV2 has 14 times fewer
                operations than ResNet-50, but it can be slower on a
                data-centre GPU.
        \end{itemize}
      \end{block}
    \end{column}
  \end{columns}
  \vskip 0.6em
  \keypoint{Use the count to compare models. Then measure the latency on
    the device (Day 9).}
  \source{\emph{Machine Learning Systems}, Volume I, chapter 6.}
\end{frame}

% -------------------------------------------------------------------- budget 3
\begin{frame}{Budget 3: activation memory and RAM}
  \begingroup\centering
  \begin{tikzpicture}[font=\footnotesize, x=1cm, y=0.065cm,
      >={Stealth[length=5pt]}]
    % y unit: 1 = 1000 values
    \draw[coolgray] (-0.2,0) -- (10.4,0);
    \fill[cardinalred!70] (0.2,0) rectangle (1.2,27.6);
    \fill[cardinalred!70] (2.4,0) rectangle (3.4,36.9);
    \fill[coolgray!45]    (4.6,0) rectangle (5.6,18.4);
    \fill[coolgray!45]    (6.8,0) rectangle (7.8,9.2);
    \fill[coolgray!45]    (9.0,0) rectangle (10.0,0.6);
    \node[above] at (0.7,27.6) {27\,648};
    \node[above] at (2.9,36.9) {36\,864};
    \node[above] at (5.1,18.4) {18\,432};
    \node[above] at (7.3,9.2)  {9216};
    \node[above] at (9.5,0.6)  {64};
    \node[below] at (0.7,0) {input};
    \node[below] at (2.9,0) {output 1};
    \node[below] at (5.1,0) {output 2};
    \node[below] at (7.3,0) {output 3};
    \node[below] at (9.5,0) {pooled};
    \draw[->] (1.3,10) -- node[above] {conv 1} (2.3,10);
    \draw[->] (3.5,10) -- node[above] {conv 2} (4.5,10);
    \draw[->] (5.7,5)  -- node[above] {conv 3} (6.7,5);
    \draw[->] (7.9,3)  -- node[above] {pool} (8.9,3);
    \node[cardinalred, anchor=west] at (4.0,40) {peak: $27\,648 + 36\,864
      = 64\,512$ values};
  \end{tikzpicture}\par\endgroup
  \vskip 0.2em
  \begin{itemize}
    \item When a layer runs, its input and its output are in the RAM at the
          same time.
    \item After the layer, the runtime uses the memory of the input again.
    \item The peak activation memory is the largest sum over all layers.
  \end{itemize}
\end{frame}

\begin{frame}{The same network in two data types}
  \renewcommand{\arraystretch}{1.3}
  \begin{tabular}{@{}lrr@{}}
    \toprule
    & \textbf{float32} & \textbf{int8} \\
    \midrule
    Bytes for each value                 & 4 & 1 \\
    Model size (24\,234 parameters)      & 95 KB & 24 KB \\
    Peak activation memory (64\,512 values) & 252 KB & 63 KB \\
    \bottomrule
  \end{tabular}
  \vskip 0.8em
  A microcontroller with 256 KB of RAM:
  \begin{itemize}
    \item \code{float32}: the tensors need 252 KB. No RAM remains for the
          program, so the model does not fit.
    \item \code{int8}: the tensors need 63 KB. The model fits.
  \end{itemize}
  \vskip 0.4em
  \keypoint{This model is small in flash in the two data types. The RAM
    decides. You cannot see this from the number of parameters.}
\end{frame}

\begin{frame}{The budgets are independent}
  \small
  \renewcommand{\arraystretch}{1.3}
  \begin{tabular}{@{}L{0.36\textwidth}rrr@{}}
    \toprule
    \textbf{One layer} & \textbf{Parameters} & \textbf{Model size}
      & \textbf{Output tensor} \\
    \midrule
    Fully connected, 1024 to 1024 & 1\,049\,600 & 4.0 MB & 4 KB \\
    Convolution $3 \times 3$, 3 to 16 channels, image $224 \times 224$
      & 448 & 2 KB & 3.1 MB \\
    \bottomrule
  \end{tabular}
  \vskip 0.3em
  {\footnotesize\color{coolgray} All sizes are for \code{float32}.\par}
  \vskip 0.6em
  \normalsize
  \begin{itemize}
    \item The first layer needs much flash and almost no RAM.
    \item The second layer needs almost no flash and much RAM.
    \item A pruned model has fewer parameters. Its tensors can have the
          same size as before.
  \end{itemize}
  \keypoint{Report three numbers for each model: parameters, operations,
    and peak activation memory.}
\end{frame}

% ------------------------------------------------------------ efficient blocks
\begin{frame}{Efficient block 1: depthwise separable convolution}
  \begingroup\centering
  \begin{tikzpicture}[font=\footnotesize]
    % ---- standard convolution
    \node[anchor=west, font=\small\bfseries] at (-0.1,5.0)
      {Standard convolution: one step};
    \cellgrid{0.26}{3.36}{6}{0.22}{darkcyan!55}
    \cellgrid{0.13}{3.23}{6}{0.22}{treegreen!55}
    \cellgrid{0}{3.1}{6}{0.22}{cardinalred!55}
    \node at (2.2,3.85) {\large$\ast$};
    \cellgrid{3.06}{3.71}{3}{0.22}{orange!45}
    \cellgrid{2.93}{3.58}{3}{0.22}{orange!45}
    \cellgrid{2.8}{3.45}{3}{0.22}{orange!45}
    \node[anchor=west, align=left] at (4.1,3.85)
      {$C_{\text{out}}$ filters of size\\$K \times K \times C_{\text{in}}$};
    \node at (7.7,3.85) {\large$=$};
    \foreach \o in {0.4,0.3,0.2,0.1,0} {
      \cellgrid{8.2+\o}{3.3+\o}{4}{0.22}{white}
    }
    % ---- separable convolution
    \node[anchor=west, font=\small\bfseries] at (-0.1,2.6)
      {Depthwise separable convolution: two steps};
    \cellgrid{0.26}{0.96}{6}{0.22}{darkcyan!55}
    \cellgrid{0.13}{0.83}{6}{0.22}{treegreen!55}
    \cellgrid{0}{0.7}{6}{0.22}{cardinalred!55}
    \node at (2.2,1.45) {\large$\ast$};
    \cellgrid{3.4}{1.55}{3}{0.22}{darkcyan!30}
    \cellgrid{3.1}{1.25}{3}{0.22}{treegreen!30}
    \cellgrid{2.8}{0.95}{3}{0.22}{cardinalred!30}
    \node[align=center] at (3.4,0.25) {depthwise: $C_{\text{in}}$ filters\\
      of size $K \times K \times 1$};
    \node at (5.0,1.45) {\large$\ast$};
    \cellgrid{6.16}{1.51}{1}{0.22}{violet!45}
    \cellgrid{6.08}{1.43}{1}{0.22}{violet!45}
    \cellgrid{6.0}{1.35}{1}{0.22}{violet!45}
    \node[align=center] at (6.2,0.25) {pointwise: $C_{\text{out}}$ filters\\
      of size $1 \times 1 \times C_{\text{in}}$};
    \node at (7.7,1.45) {\large$=$};
    \foreach \o in {0.4,0.3,0.2,0.1,0} {
      \cellgrid{8.2+\o}{0.9+\o}{4}{0.22}{white}
    }
  \end{tikzpicture}\par\endgroup
  \vskip 0.2em
  \small
  The depthwise step filters the space. The pointwise step mixes the
  channels. The output has the same size in the two cases.
\end{frame}

\begin{frame}{Depthwise separable convolution: the cost}
  \renewcommand{\arraystretch}{1.35}
  \begin{tabular}{@{}lll@{}}
    \toprule
    & \textbf{Weights} & \textbf{$3 \times 3$, 32 to 64 channels} \\
    \midrule
    Standard convolution & $K^2 \, C_{\text{in}} \, C_{\text{out}}$
      & 18\,432 \\
    Depthwise step       & $K^2 \, C_{\text{in}}$ & 288 \\
    Pointwise step       & $C_{\text{in}} \, C_{\text{out}}$ & 2048 \\
    Separable, total     & $K^2 \, C_{\text{in}} + C_{\text{in}} \,
      C_{\text{out}}$ & 2336 \\
    \bottomrule
  \end{tabular}
  \vskip 0.6em
  \[
    \frac{\text{separable}}{\text{standard}}
      = \frac{1}{C_{\text{out}}} + \frac{1}{K^2}
      = \frac{1}{64} + \frac{1}{9} \approx \frac{1}{7.9}
  \]
  \begin{itemize}
    \item The operations decrease by the same factor, because the output
          size is the same.
    \item The output tensor is the same. The RAM does not decrease.
  \end{itemize}
  \source{The book gives a factor of 8 to 9 for $3 \times 3$ kernels:
    \emph{Machine Learning Systems}, Volume I, chapter 6.}
\end{frame}

\begin{frame}{Efficient block 2: the inverted residual}
  \begingroup\centering
  \begin{tikzpicture}[font=\footnotesize, >={Stealth[length=5pt]},
      narrow/.style={draw=darkcyan, line width=0.8pt, fill=boxgray,
                     minimum width=0.7cm, minimum height=1.0cm},
      wide/.style={draw=cardinalred, line width=0.8pt, fill=boxgray,
                   minimum width=0.7cm, minimum height=3.0cm},
      op/.style={align=center, text=coolgray}]
    \node[narrow] (i) at (0,0) {};
    \node[wide]   (e) at (2.9,0) {};
    \node[wide]   (d) at (5.8,0) {};
    \node[narrow] (o) at (8.7,0) {};
    \draw[->] (i) -- node[op, above] {expand\\$1 \times 1$, ReLU6} (e);
    \draw[->] (e) -- node[op, above] {depthwise\\$3 \times 3$, ReLU6} (d);
    \draw[->] (d) -- node[op, above] {project\\$1 \times 1$, linear} (o);
    \draw[->, darkcyan, line width=0.8pt, rounded corners=4pt]
      (i.south) -- (0,-2.0) -- node[below] {skip connection (add)}
      (8.7,-2.0) -- (o.south);
    \node[below, darkcyan] at (0,1.1) {};
    \node[above] at (0,0.55)   {$C$};
    \node[above] at (2.9,1.55) {$t \times C$};
    \node[above] at (5.8,1.55) {$t \times C$};
    \node[above] at (8.7,0.55) {$C$};
  \end{tikzpicture}\par\endgroup
  \vskip 0.2em
  \begin{itemize}
    \item The block of MobileNetV2. The expansion factor $t$ is usually 6.
    \item The $3 \times 3$ filter operates on many channels, but it is a
          depthwise filter, so it is cheap.
    \item The skip connection joins the two narrow tensors. A classic
          residual block joins the wide tensors. Thus the name
          ``inverted''.
  \end{itemize}
  \source{Sandler et al., ``MobileNetV2'', 2018.}
\end{frame}

\begin{frame}{Inverted residual: cheap in flash, expensive in RAM}
  Example: $C = 32$, $t = 6$, feature map $28 \times 28$, stride 1.
  \vskip 0.4em
  \begin{columns}[T]
    \begin{column}{0.47\textwidth}
      \begin{block}{Weights}
        \begin{itemize}
          \item Expand: $32 \times 192 = 6144$
          \item Depthwise: $9 \times 192 = 1728$
          \item Project: $192 \times 32 = 6144$
          \item Total: 14\,016
        \end{itemize}
        A standard $3 \times 3$ convolution on 192 channels: 331\,776
        weights, 24 times more.
      \end{block}
    \end{column}
    \begin{column}{0.47\textwidth}
      \begin{block}{Tensors}
        \begin{itemize}
          \item Input: $28 \times 28 \times 32 = 25\,088$ values
          \item Expanded: 150\,528 values
          \item During the depthwise step: two expanded tensors and the
                input for the skip connection
          \item Peak: 326\,144 values. This is 13 times the input.
        \end{itemize}
      \end{block}
    \end{column}
  \end{columns}
  \vskip 0.5em
  \keypoint{The block saves parameters and operations. It does not save
    RAM. On a microcontroller, the expanded tensors are often the limit.}
  \source{Calculated for this course, with separate buffers for the input
    and the output of each layer.}
\end{frame}

\begin{frame}{The result: two models for the same task}
  \begin{columns}[T]
    \begin{column}{0.52\textwidth}
      \small
      \renewcommand{\arraystretch}{1.3}
      \setlength{\tabcolsep}{4pt}
      \begin{tabular}{@{}lrr@{}}
        \toprule
        & \textbf{ResNet-50} & \textbf{MobileNetV2} \\
        \midrule
        Parameters      & 25.6 M & 3.5 M \\
        Operations      & 4.1 G  & 300 M \\
        Size, float32   & 98 MB  & 13.4 MB \\
        \bottomrule
      \end{tabular}
      \vskip 0.6em
      \normalsize
      Both models classify ImageNet images.
    \end{column}
    \begin{column}{0.44\textwidth}
      \begin{itemize}
        \item 7 times fewer parameters.
        \item 14 times fewer operations.
        \item The cause is the block design, not a smaller input image.
        \item MobileNetV2 still does not fit a microcontroller with 256 KB
              of RAM. Part 3 shows why.
      \end{itemize}
    \end{column}
  \end{columns}
  \vskip 0.6em
  \keypoint{The design of the blocks decides the cost before you apply
    quantization or pruning.}
  \source{Parameters and operations: \emph{Machine Learning Systems}, Volume
    I, chapter 6. Sizes with 1 MB = 1024 KB.}
\end{frame}

% ---------------------------------------------------------------------- method
\begin{frame}{Method: calculate the cost of a model}
  \begin{enumerate}
    \item Write the output size of each layer: $H_{\text{out}} \times
          W_{\text{out}} \times C_{\text{out}}$.
    \item Calculate the parameters of each layer. Add them.\\
          $\rightarrow$ model size $\rightarrow$ compare with the flash.
    \item Calculate the MACs of each layer. Add them.\\
          $\rightarrow$ first estimate of the latency and of the energy.
    \item For each layer, add the size of its input and of its output.
          Take the maximum.\\
          $\rightarrow$ peak activation memory $\rightarrow$ compare with
          the RAM.
  \end{enumerate}
  \vskip 0.6em
  \begin{block}{In the lab}
    A notebook does these four steps for three models. You compare the
    results with the budgets of four devices.
  \end{block}
\end{frame}

% -------------------------------------------------------------------- exercise
\exerciseframe{three layers by hand}{%
  Calculate the output size and the number of parameters of each layer.
  Include the biases.
  \begin{enumerate}
    \item \textbf{Convolution.} Input $96 \times 96 \times 3$. 16 filters of
          size $3 \times 3$, stride 2, padding 1.
    \item \textbf{Depthwise separable convolution.} Input
          $48 \times 48 \times 16$. Depthwise step: $3 \times 3$, stride 1,
          padding 1. Pointwise step: 32 filters.
    \item \textbf{Fully connected layer.} Input: 32 values from a global
          average pooling layer. Output: 10 classes.
  \end{enumerate}
  Then answer: how many parameters does a standard $3 \times 3$ convolution
  need in place of layer 2?
  \vskip 0.3em
  Work alone for 7 minutes. Use no calculator.
}

\answerframe{three layers by hand}{%
  \begin{enumerate}
    \item Output: $\lfloor (96 + 2 - 3) / 2 \rfloor + 1 = 48$, so
          $48 \times 48 \times 16$.\\
          Parameters: $3 \times 3 \times 3 \times 16 + 16 = 448$.
    \item Output: $48 \times 48 \times 32$. The stride is 1 and the padding
          keeps the size.\\
          Depthwise: $3 \times 3 \times 16 + 16 = 160$.
          Pointwise: $16 \times 32 + 32 = 544$.\\
          Parameters: $160 + 544 = 704$.
    \item Output: 10 values. Parameters: $32 \times 10 + 10 = 330$.
  \end{enumerate}
  Standard convolution in place of layer 2:
  $3 \times 3 \times 16 \times 32 + 32 = 4640$ parameters. The separable
  layer needs 6.6 times fewer.
}
